\documentclass[10pt,a4paper]{amsart}
\usepackage{amsmath,amssymb,amsthm,graphicx}
% Retained minimal preamble for the Stacks Project chapter more-algebra.tex (More on Algebra).
% The chapter file pinned-original-source.tex carries only \input{preamble} on its first line, so no
% preamble text lives inside the chapter itself.  The Stacks Project uses one single shared preamble for
% every chapter; the snapshot stages it identically as sources/stacks-chow/preamble.tex and
% sources/stacks-etale-cohomology/preamble.tex.  The block below is copied mechanically, line for line,
% from sources/stacks-chow/preamble.tex lines 14-18 (the Xy-pic setup) and 154-217 (the native theorem
% declarations together with the mathematical macros).  Nothing was added, removed or redefined; the
% wrapper class and the amsmath/amssymb/amsthm/graphicx base packages are supplied by the assembler.
% Chapter: more-algebra.tex -- More on Algebra.
% For commutative diagrams we use Xy-pic
\usepackage[all]{xy}

% We use 2cell for 2-commutative diagrams.
\xyoption{2cell}

% Theorem environments.
%
\theoremstyle{plain}
\newtheorem{theorem}[subsection]{Theorem}
\newtheorem{proposition}[subsection]{Proposition}
\newtheorem{lemma}[subsection]{Lemma}

\theoremstyle{definition}
\newtheorem{definition}[subsection]{Definition}
\newtheorem{example}[subsection]{Example}
\newtheorem{exercise}[subsection]{Exercise}
\newtheorem{situation}[subsection]{Situation}

\theoremstyle{remark}
\newtheorem{remark}[subsection]{Remark}
\newtheorem{remarks}[subsection]{Remarks}

\numberwithin{equation}{subsection}

% Macros
%
\def\lim{\mathop{\mathrm{lim}}\nolimits}
\def\colim{\mathop{\mathrm{colim}}\nolimits}
\def\Spec{\mathop{\mathrm{Spec}}}
\def\Hom{\mathop{\mathrm{Hom}}\nolimits}
\def\Ext{\mathop{\mathrm{Ext}}\nolimits}
\def\SheafHom{\mathop{\mathcal{H}\!\mathit{om}}\nolimits}
\def\SheafExt{\mathop{\mathcal{E}\!\mathit{xt}}\nolimits}
\def\Sch{\mathit{Sch}}
\def\Mor{\mathop{\mathrm{Mor}}\nolimits}
\def\Ob{\mathop{\mathrm{Ob}}\nolimits}
\def\Sh{\mathop{\mathit{Sh}}\nolimits}
\def\NL{\mathop{N\!L}\nolimits}
\def\CH{\mathop{\mathrm{CH}}\nolimits}
\def\proetale{{pro\text{-}\acute{e}tale}}
\def\etale{{\acute{e}tale}}
\def\QCoh{\mathit{QCoh}}
\def\Ker{\mathop{\mathrm{Ker}}}
\def\Im{\mathop{\mathrm{Im}}}
\def\Coker{\mathop{\mathrm{Coker}}}
\def\Coim{\mathop{\mathrm{Coim}}}

% Boxtimes
%
\DeclareMathSymbol{\boxtimes}{\mathbin}{AMSa}{"02}

%
% Macros for moduli stacks/spaces
%
\def\QCohstack{\mathcal{QC}\!\mathit{oh}}
\def\Cohstack{\mathcal{C}\!\mathit{oh}}
\def\Spacesstack{\mathcal{S}\!\mathit{paces}}
\def\Quotfunctor{\mathrm{Quot}}
\def\Hilbfunctor{\mathrm{Hilb}}
\def\Curvesstack{\mathcal{C}\!\mathit{urves}}
\def\Polarizedstack{\mathcal{P}\!\mathit{olarized}}
\def\Complexesstack{\mathcal{C}\!\mathit{omplexes}}
% \Pic is the operator that assigns to X its picard group, usage \Pic(X)
% \Picardstack_{X/B} denotes the Picard stack of X over B
% \Picardfunctor_{X/B} denotes the Picard functor of X over B
\def\Pic{\mathop{\mathrm{Pic}}\nolimits}
\def\Picardstack{\mathcal{P}\!\mathit{ic}}
\def\Picardfunctor{\mathrm{Pic}}
\def\Deformationcategory{\mathcal{D}\!\mathit{ef}}

% --- appended 2026-10-12 (Chief): the project own definitions of the reference/slogan/history
% display environments, copied mechanically from the shared Stacks preamble lines 8-12
% (\usepackage{verbatim} plus the three \newenvironment lines, which the chapter preamble keeps
% verbatim from the project preamble). Nothing added, removed or redefined: these environments
% carry no mathematical content and the project itself typesets them as comments. Without these
% three lines a counted statement whose own text uses them fails with Environment slogan undefined.
\usepackage{verbatim}
\newenvironment{reference}{\comment}{\endcomment}
\newenvironment{slogan}{\comment}{\endcomment}
\newenvironment{history}{\comment}{\endcomment}
\usepackage{cleveref}
\crefname{originallemma}{original lemma}{original lemmas}
\Crefname{originallemma}{Original lemma}{Original lemmas}
\makeatletter
\newlabel{algebra-lemma-module-colimit-fp}{{00HA}{source}{}{sourcepointer-1}{}}
\newlabel{algebra-lemma-module-colimit-fp@cref}{{[originallemma][0][]00HA}{[1][0][]source}{}{}{}}
\makeatother
\crefname{originallemma}{original lemma}{original lemmas}
\Crefname{originallemma}{Original lemma}{Original lemmas}
\makeatletter
\newlabel{lemma-perfect}{{18889}{source}{}{sourcepointer-2}{}}
\newlabel{lemma-perfect@cref}{{[originallemma][0][]18889}{[1][0][]source}{}{}{}}
\makeatother
\crefname{originallemma}{original lemma}{original lemmas}
\Crefname{originallemma}{Original lemma}{Original lemmas}
\makeatletter
\newlabel{lemma-projective-amplitude}{{17414}{source}{}{sourcepointer-3}{}}
\newlabel{lemma-projective-amplitude@cref}{{[originallemma][0][]17414}{[1][0][]source}{}{}{}}
\makeatother
\crefname{originallemma}{original lemma}{original lemmas}
\Crefname{originallemma}{Original lemma}{Original lemmas}
\makeatletter
\newlabel{lemma-pseudo-coherence-colimit-ext}{{16596}{source}{}{sourcepointer-4}{}}
\newlabel{lemma-pseudo-coherence-colimit-ext@cref}{{[originallemma][0][]16596}{[1][0][]source}{}{}{}}
\makeatother
\crefname{originallemma}{original lemma}{original lemmas}
\Crefname{originallemma}{Original lemma}{Original lemmas}
\makeatletter
\newlabel{lemma-tor-amplitude}{{16906}{source}{}{sourcepointer-5}{}}
\newlabel{lemma-tor-amplitude@cref}{{[originallemma][0][]16906}{[1][0][]source}{}{}{}}
\makeatother
\title{Projective amplitude of a pseudo-coherent complex is equivalent to perfection with bounded tor amplitude and to Ext vanishing against finitely presented modules}
\author{The Stacks Project Authors}
\date{stacks-project:more-algebra.tex:local-snapshot-20261012; GNU FDL 1.2 or later}
\begin{document}
\maketitle
\setcounter{subsection}{1}
\noindent\textit{Editorial note.} All delivered mathematical text is the author's own, copied from the pinned Stacks Project chapter file more-algebra.tex; no statement or proof step was rewritten, reordered or removed.\par\medskip
\noindent\textit{Editorial note.} The chapter file pinned-original-source.tex carries only \textbackslash{}input\{preamble\} on its first line, so no preamble text lives inside the chapter. The retained file sources/stacks-more-algebra/retained-preamble.tex is a mechanical, line-for-line copy of the project's single shared preamble as staged in the snapshot; nothing was added, removed or redefined and no package was added.\par\medskip
\noindent\textit{Editorial note.} The source URL is the source-level Stacks Project URL, as required when more than one card is produced from one source; the per-claim tag is carried by claim\_key.\par\medskip
\noindent\textit{Editorial note.} The author writes 'the assumed vanishing of ext groups' with a lowercase 'ext' in the counted proof span; that is the author's text and is reproduced unchanged. The six references inside the counted proof span are recorded as non-load-bearing navigation pointers.\par\medskip
\noindent\textit{Allowed standard prerequisites (editorial).}\begin{itemize}
\item The characterisation of perfect complexes and the relation between projective amplitude and tor amplitude. invoked at proof lines 20153-20158. Establishes the equivalence of the first two conditions.
\item Vanishing of Ext against finitely presented modules versus against arbitrary modules, for a pseudo-coherent K. used at proof lines 20161-20165. Is the reduction that brings conditions (3), (4) and (5) into the range of the projective-amplitude criterion.
\item The projective-amplitude criterion in terms of Ext vanishing. invoked at proof lines 20166-20167. Gives the equivalence of (1) with the three numerical conditions.
\end{itemize}\medskip
\ifdefined\hypertarget\hypertarget{sourcepointer-1}{}\fi
\noindent\textit{Source reference.} Original lemma 00HA: focused-seeds/AGP/sources/stacks-snapshot-20261011/tags.txt, source lines 636--636. Cross-chapter navigation pointer into the Stacks Project; the standard fact it supplies is applied inside the counted proof span and the argument resting on it stays there.\par\medskip
\ifdefined\hypertarget\hypertarget{sourcepointer-2}{}\fi
\noindent\textit{Source reference.} Original lemma 18889: sources/stacks-more-algebra/pinned-original-source.tex, source lines 18889--18889. Same-chapter navigation pointer; the statement it names is standard for the section and the argument that consumes it stays inside the counted proof span.\par\medskip
\ifdefined\hypertarget\hypertarget{sourcepointer-3}{}\fi
\noindent\textit{Source reference.} Original lemma 17414: sources/stacks-more-algebra/pinned-original-source.tex, source lines 17414--17414. Same-chapter navigation pointer; the statement it names is standard for the section and the argument that consumes it stays inside the counted proof span.\par\medskip
\ifdefined\hypertarget\hypertarget{sourcepointer-4}{}\fi
\noindent\textit{Source reference.} Original lemma 16596: sources/stacks-more-algebra/pinned-original-source.tex, source lines 16596--16596. Same-chapter navigation pointer; the statement it names is standard for the section and the argument that consumes it stays inside the counted proof span.\par\medskip
\ifdefined\hypertarget\hypertarget{sourcepointer-5}{}\fi
\noindent\textit{Source reference.} Original lemma 16906: sources/stacks-more-algebra/pinned-original-source.tex, source lines 16906--16906. Same-chapter navigation pointer; the statement it names is standard for the section and the argument that consumes it stays inside the counted proof span.\par\medskip
\begin{lemma}
\label{lemma-projective-amplitude-pseudo-coherent}
Let $R$ be a ring. Let $K$ be a pseudo-coherent object of $D(R)$.
Let $a, b \in \mathbf{Z}$. The following are equivalent
\begin{enumerate}
\item $K$ has projective-amplitude in $[a, b]$,
\item $K$ is perfect of tor-amplitude in $[a, b]$,
\item $\Ext^i_R(K, N) = 0$ for all finitely presented $R$-modules $N$ and all
$i \not \in [-b, -a]$,
\item $H^n(K) = 0$ for $n > b$ and
$\Ext^i_R(K, N) = 0$ for all finitely presented $R$-modules $N$ and
all $i > -a$, and
\item $H^n(K) = 0$ for $n \not \in [a - 1, b]$ and
$\Ext^{-a + 1}_R(K, N) = 0$ for all finitely presented $R$-modules $N$.
\end{enumerate}
\end{lemma}

\begin{proof}
From the final statement of Lemma \ref{lemma-perfect} we see that
(2) implies (1).  If (1) holds, then $K$ can be represented by
a complex of projective modules $P^i$ with $P^i = 0$ for $i \not \in [a, b]$.
Since projective modules are flat (as summands of free modules), we see
that $K$ has tor-amplitude in $[a, b]$, see Lemma \ref{lemma-tor-amplitude}.
Thus by Lemma \ref{lemma-perfect} we see that (2) holds.

\medskip\noindent
In conditions (3), (4), (5) the assumed vanishing of ext groups
$\Ext^i_R(K, M)$ for $M$ of finite presentation is equivalent to the
vanishing for all $R$-modules $M$ by
Lemma \ref{lemma-pseudo-coherence-colimit-ext}
and Algebra, Lemma \ref{algebra-lemma-module-colimit-fp}.
Thus the equivalence of (1), (3), (4), and (5)
follows from Lemma \ref{lemma-projective-amplitude}.
\end{proof}

\end{document}
